\documentclass[14pt,a4paper,oneside]{extbook}
\usepackage[utf8]{vietnam}
\usepackage{amsmath,amssymb,mathrsfs,fancyhdr,enumerate,multirow,makecell}
\usepackage{tkz-euclide,tikz-3dplot,tikz,tkz-tab}
\usetikzlibrary{arrows.meta,calc,intersections,angles,quotes,backgrounds,shapes.geometric,patterns,through,shadows,positioning}
\usepackage{yhmath}
\usepackage[top=0.7cm, bottom=0.7cm, left=1.0cm, right=1.0cm, headheight=16pt, headsep=3mm, footskip=5mm]{geometry}
\usepackage[hidelinks,unicode]{hyperref}
\usepackage{qrcode}
\usepackage[many]{tcolorbox}
\usepackage{fontawesome}

\pagestyle{fancy}
\fancyhf{}
\fancyhead[L]{\footnotesize\bfseries\color{teal!85!black}\faBook~SỔ TAY KIẾN THỨC TOÁN 5}
\fancyhead[R]{\footnotesize\bfseries\color{blue!85!black}Lớp Toán Thầy Lĩnh -- loptoanthaylinh.vn}
\fancyfoot[C]{\footnotesize\bfseries\color{gray!80!black}-- Sổ tay kiến thức: Bài 30. Ôn tập số thập phân --}
\renewcommand{\headrulewidth}{0.4pt}

\begin{document}

% --- BANNER TIÊU ĐỀ BÀI HỌC ---
\noindent
\begin{tcolorbox}[
  enhanced,
  colback=teal!8!white,
  colframe=teal!80!black,
  arc=3mm,
  boxrule=1.1pt,
  width=\linewidth,
  halign=center,
  top=2.2mm, bottom=2.2mm,
  left=3mm, right=3mm,
  drop fuzzy shadow
]
  {\large\bfseries\color{teal!95!black}\faBookmark\ BÀI 30. ÔN TẬP SỐ THẬP PHÂN (ÔN TẬP HỌC KÌ 1)}\\[2pt]
  {\scriptsize\color{black!75}\itshape Tổng ôn cấu tạo, các hàng, so sánh và làm tròn số thập phân chuẩn bị thi Học kì 1 trên 1 trang A4}
\end{tcolorbox}

\vspace{3.2mm}
% --- KHỐI 1: CẤU TẠO VÀ ĐỌC VIẾT STP ---
\noindent
\begin{minipage}[t]{0.56\textwidth}
  \begin{tcolorbox}[
    enhanced,
    title={\footnotesize\bfseries\color{white}\faKey\ 1. Cấu tạo và giá trị chữ số},
    colback=blue!2!white,
    colframe=blue!75!black,
    arc=2.5mm, boxrule=0.8pt, top=1.8mm, bottom=1.8mm, left=2.2mm, right=2.2mm
  ]
    \footnotesize
    $\bullet$ Số thập phân gồm: \textbf{Phần nguyên} (trước dấu phẩy) và \textbf{Phần thập phân} (sau dấu phẩy).\\
    $\bullet$ Hàng sau dấu phẩy: \textbf{Phần mười} ($\frac{1}{10}$), \textbf{phần trăm} ($\frac{1}{100}$), \textbf{phần nghìn} ($\frac{1}{1000}$)...\\
    \textit{VD:} Trong số $204,507$: Chữ số $4$ ở hàng đơn vị; chữ số $5$ có giá trị $\frac{5}{10}$; chữ số $7$ có giá trị $\frac{7}{1000}$.
  \end{tcolorbox}
\end{minipage}%
\hfill
\begin{minipage}[t]{0.42\textwidth}
  \begin{tcolorbox}[
    enhanced,
    title={\footnotesize\bfseries\color{white}\faExchange\ 2. Số thập phân bằng nhau},
    colback=teal!2!white,
    colframe=teal!75!black,
    arc=2.5mm, boxrule=0.8pt, top=1.8mm, bottom=1.8mm, left=2mm, right=2mm
  ]
    \footnotesize
    $\bullet$ Viết thêm chữ số $0$ vào tận cùng bên phải phần thập phân: giá trị \textbf{không đổi}.\\
    $0,7 = 0,70 = 0,700$.\\
    $\bullet$ Bỏ bớt chữ số $0$ ở tận cùng bên phải phần thập phân: giá trị \textbf{không đổi}.\\
    $15,400 = 15,4$.
  \end{tcolorbox}
\end{minipage}

\vspace{3mm}

% --- KHỐI 2: SO SÁNH VÀ LÀM TRÒN SỐ THẬP PHÂN ---
\noindent
\begin{tcolorbox}[
  enhanced,
  title={\footnotesize\bfseries\color{white}\faBalanceScale\ 3. So sánh và Làm tròn số thập phân},
  colback=purple!2!white,
  colframe=purple!70!black,
  arc=2.5mm, boxrule=0.8pt, top=1.8mm, bottom=1.8mm, left=2.5mm, right=2.5mm
]
  \footnotesize
  \begin{tabular}{p{0.48\textwidth}|p{0.48\textwidth}}
    \textbf{SO SÁNH HAI SỐ THẬP PHÂN} & \textbf{LÀM TRÒN SỐ THẬP PHÂN} \\ \hline
    $\bullet$ So sánh phần nguyên trước: $15,2 > 9,99$. & $\bullet$ Nhìn chữ số ngay bên phải hàng làm tròn:\\
    $\bullet$ Phần nguyên bằng nhau $\rightarrow$ so từng hàng tương ứng từ trái sang phải: $4,58 < 4,6$ (vì $5 < 6$). & $\triangleright$ Nếu $< 5$: Giữ nguyên chữ số hàng làm tròn.\\
    $\bullet$ Mẹo: Thêm số $0$ vào cho bằng số chữ số thập phân rồi so sánh: $4,58$ so với $4,60$. & $\triangleright$ Nếu $\ge 5$: Tăng thêm $1$ đơn vị.\\
    \textit{VD:} $3,74 \approx 3,7$ (hàng phần mười); $3,75 \approx 3,8$.
  \end{tabular}
\end{tcolorbox}

\vspace{3mm}

% --- KHỐI 3: SƠ ĐỒ TƯ DUY ---
\noindent
\begin{tcolorbox}[
  enhanced,
  title={\footnotesize\bfseries\color{white}\faSitemap\ Sơ đồ tư duy: Ôn tập Số thập phân HKI},
  colback=yellow!2!white,
  colframe=teal!85!black,
  arc=3mm, boxrule=0.9pt, top=1.8mm, bottom=1.8mm, left=1.5mm, right=1.5mm,
  halign=center
]
  \begin{tikzpicture}[
    scale=0.92,
    every node/.style={transform shape},
    root/.style={rectangle, rounded corners=5pt, fill=teal!85!black, text=white, font=\bfseries\footnotesize, inner sep=5pt, align=center},
    b1/.style={rectangle, rounded corners=4pt, fill=blue!12, draw=blue!75!black, text=blue!90!black, font=\bfseries\scriptsize, inner sep=4.5pt, align=center},
    b2/.style={rectangle, rounded corners=4pt, fill=purple!12, draw=purple!75!black, text=purple!90!black, font=\bfseries\scriptsize, inner sep=4.5pt, align=center},
    b3/.style={rectangle, rounded corners=4pt, fill=orange!12, draw=orange!75!black, text=orange!90!black, font=\bfseries\scriptsize, inner sep=4.5pt, align=center}
  ]
    \node[root] (core) at (0, 0.85) {\faBook\ SỐ THẬP PHÂN HKI};
    \node[b1] (n1) at (-5.2, -0.45) {\textbf{CẤU TẠO:}\\[2pt]Phần nguyên + Phần thập phân};
    \node[b2] (n2) at (0, -0.65) {\textbf{ĐỔI ĐƠN VỊ:}\\[2pt]Gắn với phân số thập phân};
    \node[b3] (n3) at (5.2, -0.45) {\textbf{SO SÁNH \& LÀM TRÒN:}\\[2pt]Theo hàng / Xét chữ số phải};

    \draw[teal!85!black, thick, -{Stealth[length=2mm]}] (core) -| (n1);
    \draw[teal!85!black, thick, -{Stealth[length=2mm]}] (core) -- (n2);
    \draw[teal!85!black, thick, -{Stealth[length=2mm]}] (core) -| (n3);
  \end{tikzpicture}
\end{tcolorbox}

\vspace{3mm}

% --- KHỐI 4: CHUYỂN ĐỔI NHANH & BẪY TRẮC NGHIỆM ---
\noindent
\begin{minipage}[t]{0.54\textwidth}
  \begin{tcolorbox}[
    enhanced,
    title={\footnotesize\bfseries\color{white}\faBolt\ Dạng toán viết số thập phân theo cấu tạo},
    colback=green!3!white,
    colframe=green!60!black,
    arc=2.5mm, boxrule=0.8pt, top=1.8mm, bottom=1.8mm, left=2.2mm, right=2.2mm
  ]
    \footnotesize
    ``Số thập phân gồm $7$ chục, $2$ đơn vị, $5$ phần trăm viết là:'':\\
    $\bullet$ Hàng chục: $7$; hàng đơn vị: $2 \rightarrow$ Phần nguyên là $72$.\\
    $\bullet$ Hàng phần mười không có $\rightarrow$ viết chữ số \textbf{0}.\\
    $\bullet$ Hàng phần trăm: $5 \rightarrow$ Kết quả: $\mathbf{72,05}$.
  \end{tcolorbox}
\end{minipage}%
\hfill
\begin{minipage}[t]{0.44\textwidth}
  \begin{tcolorbox}[
    enhanced,
    title={\footnotesize\bfseries\color{white}\faExclamationTriangle\ Mẹo tránh bẫy trắc nghiệm},
    colback=red!2!white,
    colframe=red!70!black,
    arc=2.5mm, boxrule=0.8pt, top=1.8mm, bottom=1.8mm, left=2.2mm, right=2.2mm
  ]
    \footnotesize
    $\bullet$ Quên viết chữ số $0$ ở hàng khuyết (như viết $72,5$ thay vì $72,05$).\\
    $\bullet$ Số $0$ đứng giữa các chữ số không được bỏ đi ($3,04 \neq 3,4$).
  \end{tcolorbox}
\end{minipage}

\end{document}
